\begin{figure}[ht]
\centering
\begin{tikzpicture}[x=1.4cm, y=1cm, font=\small]

% Axes: x = number of factors pinned, y = log10 of the bound value
\draw[->, thick] (0, 0) -- (5.6, 0) node[anchor=north east] {factors pinned};
\draw[->, thick] (0, 0) -- (0, 5.3) node[anchor=south east] {$\log_{10}(\text{bound})$};

% x ticks
\foreach \k in {0, 1, 2, 3, 4, 5} {
  \draw (\k, -0.05) -- (\k, 0.05);
  \node[anchor=north, font=\scriptsize] at (\k, -0.05) {\k};
}
% y ticks
\foreach \y/\lbl in {0.5/4, 1.5/5, 2.5/6, 3.5/7, 4.5/8} {
  \draw (-0.05, \y) -- (0.05, \y);
  \node[anchor=east, font=\scriptsize] at (-0.05, \y) {$10^{\lbl}$};
}

% Upper bound descending
\draw[thick, engred]
  plot coordinates {(0, 4.9) (1, 4.2) (2, 3.6) (3, 3.15) (4, 2.9) (5, 2.75)};
\node[anchor=west, engred, font=\footnotesize] at (5, 2.75) {upper bound};

% Lower bound ascending
\draw[thick, engblue]
  plot coordinates {(0, 0.4) (1, 1.0) (2, 1.5) (3, 1.9) (4, 2.15) (5, 2.3)};
\node[anchor=west, engblue, font=\footnotesize] at (5, 2.3) {lower bound};

% Shaded gap between bounds
\fill[enggray, opacity=0.15]
  plot coordinates {(0, 4.9) (1, 4.2) (2, 3.6) (3, 3.15) (4, 2.9) (5, 2.75)}
  -- plot coordinates {(5, 2.3) (4, 2.15) (3, 1.9) (2, 1.5) (1, 1.0) (0, 0.4)}
  -- cycle;

% True value line
\draw[thick, dashed, engamber] (0, 2.55) -- (5.4, 2.55);
\node[anchor=west, engamber, font=\scriptsize] at (3.2, 2.72) {true value};

% Annotation of the factor-of-three band around the close
\draw[<->, enggray] (5.15, 2.3) -- (5.15, 2.75);
\node[anchor=west, enggray, font=\scriptsize] at (5.15, 2.52) {$\approx 3\times$};

\end{tikzpicture}
\caption[Upper and lower bounds narrowing as factors are pinned]{The
two-sided bounding move under deep uncertainty. Each factor pinned to
a defensible value pulls the upper bound down and the lower bound up;
the shaded region is the range that survives. When the two bounds
close to within a factor of three of each other, the problem has
reverted to ordinary Fermi estimation and a point estimate is
reportable. When the bounds refuse to close below an order of
magnitude after every factor is pinned, the width itself is the honest
answer, and the factor whose two sides remain farthest apart is the
one the postmortem names.}
\label{fig:vol01:ch09:bounds-narrowing}
\end{figure}
